\documentclass[12pt, a4paper]{article}

% ─── Packages ────────────────────────────────────────────────────────────────
\usepackage{amsmath, amssymb, amsthm}
\usepackage{mathtools}
\usepackage{geometry}
\usepackage{hyperref}
\usepackage{booktabs}
\usepackage{array}
\usepackage{enumitem}
\usepackage{graphicx}
\usepackage{xcolor}
\usepackage{algorithm}
\usepackage{algpseudocode}
\usepackage{natbib}
\usepackage{microtype}
\usepackage{cleveref}
\usepackage{bbm}
\usepackage{bm}

\geometry{a4paper, margin=1.1in, top=1.2in, bottom=1.2in}

\hypersetup{
    colorlinks=true,
    linkcolor=blue!60!black,
    citecolor=teal!70!black,
    urlcolor=blue!60!black
}

% ─── Theorem environments ─────────────────────────────────────────────────────
\newtheorem{definition}{Definition}[section]
\newtheorem{proposition}{Proposition}[section]
\newtheorem{remark}{Remark}[section]

% ─── Notation shorthands ─────────────────────────────────────────────────────
\newcommand{\R}{\mathbb{R}}
\newcommand{\E}{\mathbb{E}}
\newcommand{\Var}{\mathrm{Var}}
\newcommand{\Cov}{\mathrm{Cov}}
\newcommand{\Normal}{\mathcal{N}}
\newcommand{\KL}{\mathrm{KL}}
\newcommand{\I}{\mathbbm{1}}
\newcommand{\mW}{\mu_{W}}
\newcommand{\sW}{\sigma^{2}_{W}}
\newcommand{\mV}{\mu_{V}}
\newcommand{\sV}{\sigma^{2}_{V}}
\newcommand{\mupost}{\mu^{+}}
\newcommand{\sigpost}{\sigma^{2,+}}
\newcommand{\target}{y}
\newcommand{\dd}{\mathrm{d}}

% ─── Title ───────────────────────────────────────────────────────────────────
\title{
    \Large\textbf{A Tractable Analytic Bayesian Spiking Neural Network}\\[0.4em]
    \large Integrating TAGI and TAGI-V for Principled Inference,\\
    Homeostasis, Structural Plasticity, and Neuron Lifecycle Management
}
\author{}
\date{March 2026}

% ─────────────────────────────────────────────────────────────────────────────
\begin{document}
\maketitle
\tableofcontents
\newpage

% =============================================================================
\section{Introduction and Motivation}
\label{sec:intro}
% =============================================================================

\subsection{The problem with standard artificial neural networks}

Modern deep learning succeeds at supervised pattern recognition but deviates sharply from
how biological neural computation works. Multi-Layer Perceptrons (MLPs) and their
convolutional relatives operate through four assumptions that are biologically implausible:
(i) neurons communicate via continuous, real-valued activations rather than discrete
spikes; (ii) the topology is fixed at design time; (iii) learning requires a global
backward pass through differentiable operations (backpropagation), which has no known
biological substrate; and (iv) all neurons are always active, ignoring the energetic
imperative of sparse coding that governs cortical computation.

These simplifications are also algorithmic liabilities. Fixed topology prevents the model
from allocating representational capacity where it is needed. Backpropagation through
time (BPTT) scales quadratically in sequence length and requires storing the full
computational graph. Non-spiking activations destroy temporal structure.

\subsection{What we want to model from the brain}

The present work is guided by five empirical principles drawn from computational
neuroscience:

\begin{enumerate}[label=\textbf{P\arabic*.}, leftmargin=2.2em]
    \item \textbf{Spike-based communication.} Cortical neurons communicate via
    all-or-nothing action potentials. Information is encoded in the timing and rate of
    spikes, not in a continuous graded signal \citep{mainen1995reliability}.

    \item \textbf{Homeostatic regulation.} Neurons actively regulate their own
    excitability to maintain a target firing rate. When a neuron fires too often its
    threshold rises; when it is silent its threshold falls. This prevents runaway
    excitation or complete silencing \citep{turrigiano1999homeostatic}.

    \item \textbf{Structural plasticity.} Synapses are not fixed. Axonal sprouting and
    dendritic retraction continuously reshape connectivity based on activity levels.
    Neurons that are active but poorly connected grow new synapses; over-connected
    inactive neurons lose them \citep{bhatt2009dendritic}.

    \item \textbf{Predictive coding and surprise.} The brain does not respond uniformly
    to all stimuli. It maintains internal predictions, and neural responses are
    proportional to \emph{prediction error} --- the gap between expectation and
    observation. Novel or unexpected stimuli trigger strong, widespread responses;
    predictable stimuli are suppressed \citep{rao1999predictive}.

    \item \textbf{Neuron lifecycle.} Neurons are born, learn, mature, and die. Synaptic
    weights that never change carry no information; the brain eliminates such synapses
    through pruning, particularly during development and consolidation
    \citep{huttenlocher1979synaptic}.
\end{enumerate}

\subsection{Why standard SNNs are not enough}

First-generation Spiking Neural Networks (SNNs) capture principle P1 but address the
others only heuristically. Threshold adaptation is typically a hand-tuned decay rule with
no notion of statistical surprise. Structural plasticity is triggered by scalar
thresholds on synaptic magnitude. Weight learning either uses biologically motivated but
theoretically weak Hebbian rules (STDP), or requires surrogate gradient methods that
reintroduce the differentiability assumption \citep{pfeiffer2018deep}. Temporal depth
($T_{\max}$) is a global hyperparameter chosen by the practitioner.

\subsection{Our approach: TAGI and TAGI-V as the inference engine}

We replace the heuristic components of SNNs with a principled probabilistic framework.
\emph{Tractable Approximate Gaussian Inference} (TAGI) \citep{goulet2021tractable}
propagates closed-form Gaussian distributions through the network forward pass and
performs exact posterior updates on weights given observed data, without backpropagation.
Its variance-heteroscedastic extension \emph{TAGI-V} \citep{nguyen2022tagi} additionally
learns the output variance, providing a direct, data-driven measure of prediction
uncertainty.

The key observation that drives the entire architecture is:

\begin{center}
\textit{In TAGI, the posterior variance of every weight is a continuous, monotonically
decreasing function of the amount of useful information the weight has processed.
This single quantity simultaneously encodes learning progress, neuron vitality, temporal
horizon, novelty of incoming data, and the signal-to-noise ratio of the network's
predictions.}
\end{center}

Each biological principle P1--P5 maps directly onto a mathematical object derived
from this variance trajectory:

\begin{table}[h]
\centering
\renewcommand{\arraystretch}{1.3}
\begin{tabular}{@{}lll@{}}
\toprule
\textbf{Principle} & \textbf{Biological phenomenon} & \textbf{TAGI object} \\
\midrule
P1 & Spike-based communication & Bernoulli output of threshold comparator \\
P2 & Homeostatic threshold & Prediction likelihood from TAGI-V \\
P3 & Structural plasticity & KL divergence of weight posterior from prior \\
P4 & Predictive coding / surprise & Output variance ratio $\sV^{(t)}/\sV^{(t-1)}$ \\
P5 & Neuron lifecycle / pruning & Posterior shrinkage ratio $\sW / \sigma^2_0$ \\
\bottomrule
\end{tabular}
\caption{Mapping between biological principles and TAGI mathematical objects.}
\label{tab:mapping}
\end{table}

The remainder of this document formalises each component.

% =============================================================================
\section{Network Structure and Spatial Embedding}
\label{sec:structure}
% =============================================================================

\subsection{Topology}

The network consists of $N$ neurons embedded in a bounded 2-D physical space
$\mathcal{P} \subset \R^2$. Each neuron $i$ has a fixed spatial coordinate
$\bm{p}_i \in \mathcal{P}$. Connectivity is encoded in a sparse weighted adjacency
matrix $\mathbf{W} \in \R^{N \times N}$, where $W_{ji} \neq 0$ if and only if neuron
$j$ sends a synapse to neuron $i$. The network density is defined as
$\rho = |\{(j,i) : W_{ji} \neq 0\}| / N^2$, and is maintained in the range $[0.05,
0.10]$ by structural plasticity (Section~\ref{sec:structural}).

\subsection{Probabilistic weight representation}

Unlike standard SNNs where each weight $W_{ji}$ is a point estimate, here every weight
carries a full Gaussian posterior:

\begin{equation}
    W_{ji} \sim \Normal\!\left(\mW_{ji},\; \sW_{ji}\right)
    \label{eq:weight_dist}
\end{equation}

The prior at initialisation is $W_{ji} \sim \Normal(0, \sigma^2_0)$ with
$\sigma^2_0$ a design constant. All inference is performed analytically on these
Gaussians, requiring no gradient computation.

% =============================================================================
\section{Forward Dynamics: Gaussian Voltage Propagation}
\label{sec:forward}
% =============================================================================

\subsection{Membrane potential as a random variable}

The membrane potential of neuron $i$ at time $t$ is a random variable induced by the
uncertain weights and the binary spike inputs. The pre-synaptic input is:

\begin{equation}
    V_i^{(t)} = \lambda_i^{(t)}\, V_i^{(t-1)}\left(1 - S_i^{(t-1)}\right)
    + \sum_{j=1}^{N} W_{ji}^{(t)}\, S_j^{(t-1)}
    \label{eq:voltage}
\end{equation}

where $\lambda_i^{(t)} \in (0,1)$ is the adaptive leak (Section~\ref{sec:oscillatory})
and $S_j^{(t-1)} \in \{0,1\}$ is the spike emitted by neuron $j$ at $t-1$. Because
$S_j^{(t-1)}$ is binary, the product $W_{ji} S_j^{(t-1)}$ equals $W_{ji}$ when $j$
fired and $0$ otherwise. The sum over active pre-synaptic neurons therefore preserves
Gaussianity:

\begin{proposition}[Gaussian voltage propagation]
If $W_{ji} \sim \Normal(\mW_{ji}, \sW_{ji})$ and the $W_{ji}$ are mutually independent,
then conditional on the spike history $\{S_j^{(t-1)}\}$,
\begin{equation}
    V_i^{(t)} \sim \Normal\!\left(\mV_i^{(t)},\; \sV_i^{(t)}\right)
    \label{eq:volt_dist}
\end{equation}
with
\begin{align}
    \mV_i^{(t)} &= \lambda_i^{(t)}\, \mV_i^{(t-1)}\!\left(1-S_i^{(t-1)}\right)
        + \sum_{j} \mW_{ji}^{(t)}\, S_j^{(t-1)} \label{eq:mu_v} \\[4pt]
    \sV_i^{(t)} &= \left(\lambda_i^{(t)}\right)^2\!\sV_i^{(t-1)}\!\left(1-S_i^{(t-1)}\right)
        + \sum_{j} \sW_{ji}^{(t)}\, S_j^{(t-1)} \label{eq:sig_v}
\end{align}
\end{proposition}

\begin{proof}
The result follows immediately from the linearity of expectation and the independence of
weights. For $S_j^2 = S_j$ (idempotent for binary variables), the variance reduces to
the sum over active pre-synaptic weight variances.
\end{proof}

\subsection{Spike generation}

Neuron $i$ fires when its voltage exceeds a dynamic threshold $\theta_i^{(t)}$:

\begin{equation}
    S_i^{(t)} = H\!\left(V_i^{(t)} - \theta_i^{(t)}\right)
    \label{eq:spike}
\end{equation}

where $H$ is the Heaviside function. Under the Gaussian voltage distribution
\eqref{eq:volt_dist}, the \emph{spike probability} is:

\begin{equation}
    \boxed{
    \pi_i^{(t)} \coloneqq P\!\left(S_i^{(t)} = 1\right)
    = 1 - \Phi\!\left(\frac{\theta_i^{(t)} - \mV_i^{(t)}}{\sqrt{\sV_i^{(t)}}}\right)
    }
    \label{eq:spike_prob}
\end{equation}

where $\Phi$ is the standard normal CDF. This quantity is the central object driving
homeostasis: it measures how \emph{expected} the observed spike outcome was under the
current Gaussian belief state.

% =============================================================================
\section{TAGI Weight Inference: Replacing Backpropagation}
\label{sec:tagi}
% =============================================================================

\subsection{Why backpropagation fails for SNNs}

The Heaviside function $H$ in \eqref{eq:spike} has zero gradient almost everywhere.
Backpropagation through $S_i^{(t)}$ is undefined, forcing practitioners to use surrogate
gradients --- smooth approximations of $H$ that lack biological plausibility and
introduce additional hyperparameters. TAGI sidesteps this entirely by working in the
Gaussian belief space, not the point-estimate space.

\subsection{TAGI forward-backward update rule}

Given an observed target output $\bm{y}$ (or observed spike $S_i^{(t)}$ as a training
signal), TAGI performs a one-pass update on all weight posteriors. The update follows
the Gaussian linear model structure. For a weight $W_{ji}$ that contributed to the
prediction of output $o$ with prediction $\hat{o} \sim \Normal(\mu_o, \sigma^2_o)$ and
observed value $y_o$, the posterior update is:

\begin{align}
    \mupost_{W_{ji}} &= \mW_{ji} + k_{W_{ji},o}\,\frac{y_o - \mu_o}{\sigma^2_o} \label{eq:tagi_mean} \\[4pt]
    \sigpost_{W_{ji}} &= \sW_{ji} - \frac{k_{W_{ji},o}^2}{\sigma^2_o} \label{eq:tagi_var}
\end{align}

where the \emph{gain} is the cross-covariance:
\begin{equation}
    k_{W_{ji},o} = \Cov(W_{ji},\, \hat{o})
    = \sW_{ji}\, S_j^{(t-1)}\, \frac{\partial \mu_o}{\partial \mW_{ji}}
    \label{eq:gain}
\end{equation}

For the voltage model \eqref{eq:mu_v}, the partial derivative is $S_j^{(t-1)}$,
giving $k_{W_{ji},o} = \sW_{ji}\, S_j^{(t-1)}$. Substituting:

\begin{align}
    \mupost_{W_{ji}} &= \mW_{ji}
        + \sW_{ji}\, S_j^{(t-1)}\;\frac{y_o - \mu_o}{\sigma^2_o}
        \label{eq:tagi_w_mean}\\[4pt]
    \sigpost_{W_{ji}} &= \sW_{ji}
        - \frac{\left(\sW_{ji}\right)^2 S_j^{(t-1)}}{\sigma^2_o}
        \label{eq:tagi_w_var}
\end{align}

\begin{remark}[Biological gating]
Update \eqref{eq:tagi_w_mean} is zero if $S_j^{(t-1)} = 0$: weights are only modified
when the pre-synaptic neuron actually fired. This is a direct analogue of
\emph{Hebbian coincidence detection} --- a synapse is strengthened only when it
participated in producing the output.
\end{remark}

\begin{remark}[Variance monotonically shrinks]
From \eqref{eq:tagi_w_var}, $\sigpost_{W_{ji}} \leq \sW_{ji}$ with equality if and only
if $S_j^{(t-1)} = 0$. Every informative update compresses the posterior. This monotone
shrinkage is the foundation of the neuron lifecycle model in
Section~\ref{sec:lifecycle}.
\end{remark}

\subsection{TAGI-V: learning the output variance}

Standard TAGI treats $\sigma^2_o$ as a fixed hyperparameter. TAGI-V \citep{nguyen2022tagi}
extends the model by placing a Gaussian distribution over the log-variance, allowing
$\sV_i^{(t)}$ itself to be inferred from data. This has two consequences for our
architecture:

\begin{enumerate}
    \item The predictive distribution $\Normal(\mV_i^{(t)}, \sV_i^{(t)})$ becomes
    \emph{heteroscedastic}: the network automatically allocates higher uncertainty to
    inputs it has not seen before.

    \item The inferred $\sV_i^{(t)}$ is a direct, calibrated measure of \emph{how
    surprising} an incoming spike pattern is --- the basis of our homeostatic rule in
    Section~\ref{sec:homeostasis}.
\end{enumerate}

% =============================================================================
\section{Principled Homeostatic Plasticity via TAGI-V}
\label{sec:homeostasis}
% =============================================================================

\subsection{Biological motivation}

Homeostatic synaptic scaling \citep{turrigiano1999homeostatic} adjusts neuronal
excitability to maintain stable firing rates despite changes in network activity. The key
insight is that the brain does not regulate firing rate blindly: a neuron that fires
frequently because it is genuinely processing novel, surprising information should
\emph{not} have its threshold raised. A neuron that fires frequently because it is
responding to highly predictable, redundant input \emph{should}. Standard SNNs cannot
make this distinction. TAGI-V can.

\subsection{Surprise-gated threshold update}

Define the \emph{predictive likelihood} of the observed spike outcome:

\begin{equation}
    \pi_i^{(t)} = P\!\left(S_i^{(t)} \mid \text{history}\right)
    = \begin{cases}
        1 - \Phi\!\left(\frac{\theta_i^{(t)} - \mV_i^{(t)}}{\sqrt{\sV_i^{(t)}}}\right)
        & S_i^{(t)} = 1 \\[6pt]
        \Phi\!\left(\frac{\theta_i^{(t)} - \mV_i^{(t)}}{\sqrt{\sV_i^{(t)}}}\right)
        & S_i^{(t)} = 0
    \end{cases}
    \label{eq:likelihood}
\end{equation}

$\pi_i^{(t)} \approx 1$ means the spike outcome was expected (low surprise);
$\pi_i^{(t)} \approx 0$ means it was surprising (novel data).

The revised homeostatic rule is:

\begin{equation}
    \boxed{
    \theta_i^{(t+1)} = \theta_i^{(t)}
    + \alpha_\theta\,\bigl(1 - \pi_i^{(t)}\bigr)\, S_i^{(t)}
    - \beta_\theta\,\pi_i^{(t)}\,\bigl(1 - S_i^{(t)}\bigr)
    }
    \label{eq:homeostasis}
\end{equation}

\begin{table}[h]
\centering
\renewcommand{\arraystretch}{1.25}
\begin{tabular}{@{}cccc@{}}
\toprule
$S_i^{(t)}$ & $\pi_i^{(t)}$ & $\Delta\theta_i$ & \textbf{Interpretation} \\
\midrule
1 & $\approx 0$ & $\approx +\alpha_\theta$ & Surprising spike: raise threshold slightly \\
1 & $\approx 1$ & $\approx 0$             & Expected spike: no threshold change \\
0 & $\approx 0$ & $\approx 0$             & Expected silence: no change \\
0 & $\approx 1$ & $\approx -\beta_\theta$ & Expected spike didn't occur: lower threshold \\
\bottomrule
\end{tabular}
\caption{Behaviour of the surprise-gated homeostatic rule \eqref{eq:homeostasis} across
all spike/likelihood combinations.}
\label{tab:homeostasis}
\end{table}

Compared to the original fixed $\pm\alpha_\theta/\beta_\theta$ schedule, rule
\eqref{eq:homeostasis} is \emph{information-theoretically grounded}: the threshold moves
in proportion to the neuron's prediction error, not merely its activity level.

% =============================================================================
\section{Adaptive Temporal Horizon via Variance Stationarity}
\label{sec:tmax}
% =============================================================================

\subsection{The problem of choosing $T_{\max}$}

How far back in time should the network integrate spike history? In standard SNNs
$T_{\max}$ is a fixed global hyperparameter. But different neurons encode information at
different timescales: a neuron sensitive to fast transients should have a short memory,
while a neuron integrating slow trends needs a long one. A single global $T_{\max}$ is
always wrong for at least some neurons.

\subsection{Temporal Riccati dynamics}

In TAGI's temporal propagation, the voltage variance satisfies a discrete Riccati
recursion. Under stationary input statistics and constant leak $\lambda$:

\begin{equation}
    \sV_i^{(t)} = \lambda^2\, \sV_i^{(t-1)} + \sigma^2_{\text{process}}
    \label{eq:riccati}
\end{equation}

where $\sigma^2_{\text{process}} = \sum_j \sW_{ji} \bar{S}_j$ sums the weight variances
of active synapses ($\bar{S}_j$ is the average firing rate of neuron $j$). This
recursion has a fixed point:

\begin{equation}
    \sV_{i,\infty} = \frac{\sigma^2_{\text{process}}}{1 - \lambda^2}
    \label{eq:fixed_point}
\end{equation}

When $\sV_i^{(t)}$ converges to $\sV_{i,\infty}$, the neuron has exhausted its useful
temporal integration window: additional history contributes negligible new information.

\subsection{Per-neuron $T_{\max}$ detection}

\begin{definition}[Temporal convergence]
The \emph{adaptive temporal horizon} of neuron $i$ is:
\begin{equation}
    T_{\max,i} = \min\!\left\{t \geq 1 : \left|\sV_i^{(t)} - \sV_i^{(t-1)}\right| < \varepsilon_T \right\}
    \label{eq:tmax}
\end{equation}
for a tolerance $\varepsilon_T > 0$.
\end{definition}

Beyond $T_{\max,i}$, no further temporal context is integrated for neuron $i$. This
gives each neuron its own biologically plausible time constant, analogous to the
membrane time constant $\tau_m$ in conductance-based neuron models \citep{dayan2001theoretical}.

\begin{remark}[Connection to the Kalman filter]
Equation \eqref{eq:riccati} is precisely the prediction-step variance recursion of a
Kalman filter. The fixed-point condition \eqref{eq:fixed_point} is the steady-state
Kalman gain. The network is thus equivalent to a bank of per-neuron Kalman filters whose
observation noise is set by the weight uncertainty of their inputs.
\end{remark}

% =============================================================================
\section{Neuron Lifecycle via Posterior Variance Shrinkage}
\label{sec:lifecycle}
% =============================================================================

\subsection{The key observation}

From Remark 2 in Section~\ref{sec:tagi}, $\sigpost_{W_{ji}} \leq \sW_{ji}$ after every
informative update. The posterior variance of each weight therefore records a
\emph{cumulative history of useful engagement}: a weight that has processed many
informative spikes will have $\sW_{ji} \ll \sigma^2_0$, while a weight that has never
been usefully engaged will retain $\sW_{ji} \approx \sigma^2_0$.

\subsection{Vitality score}

\begin{definition}[Neuron vitality]
The \emph{vitality} of neuron $i$ at time $t$ is:
\begin{equation}
    \rho_i^{(t)} = \frac{1}{|\mathcal{J}_i|}\sum_{j \in \mathcal{J}_i}
    \frac{\sW_{ji}^{(t)}}{\sigma^2_0}
    \label{eq:vitality}
\end{equation}
where $\mathcal{J}_i = \{j : W_{ji} \neq 0\}$ is the set of pre-synaptic neurons.
$\rho_i^{(t)} \in (0,1]$, with $\rho_i = 1$ meaning no shrinkage has occurred and
$\rho_i \approx 0$ meaning weights have been fully compressed.
\end{definition}

\begin{definition}[Instantaneous learning rate]
The \emph{learning activity} of neuron $i$ at time $t$ is:
\begin{equation}
    \Delta\rho_i^{(t)} = \left|\rho_i^{(t)} - \rho_i^{(t-1)}\right|
    \label{eq:drho}
\end{equation}
\end{definition}

\subsection{Three-state lifecycle}

Based on $\rho_i^{(t)}$ and $\Delta\rho_i^{(t)}$, each neuron is assigned a lifecycle
state at every timestep:

\begin{equation}
    \text{State}_i^{(t)} =
    \begin{cases}
        \textsc{dead}   & \rho_i^{(t)} \geq \tau_{\text{kill}} \\[4pt]
        \textsc{active} & \Delta\rho_i^{(t)} > \tau_{\text{active}} \\[4pt]
        \textsc{mature} & \rho_i^{(t)} < \tau_{\text{sat}}
                          \;\text{ and }\;
                          \Delta\rho_i^{(t)} \leq \tau_{\text{active}}
    \end{cases}
    \label{eq:lifecycle}
\end{equation}

with $\tau_{\text{kill}} > \tau_{\text{sat}}$ and the states evaluated in order.

\begin{table}[h]
\centering
\renewcommand{\arraystretch}{1.25}
\begin{tabular}{@{}llll@{}}
\toprule
\textbf{State} & \textbf{Meaning} & \textbf{Action} & \textbf{Brain analogue} \\
\midrule
\textsc{dead}   & Never engaged & Remove neuron and all edges & Synaptic pruning \\
\textsc{active} & Currently learning & Protect; candidate for growth & Critical period \\
\textsc{mature} & Role learned & Freeze weights; disable plasticity & Consolidated memory \\
\bottomrule
\end{tabular}
\caption{Neuron lifecycle states, actions, and biological analogues.}
\label{tab:lifecycle}
\end{table}

\begin{remark}[Safety of dead-neuron removal]
Killing a \textsc{dead} neuron is provably safe. $\rho_i \approx 1$ means
$\sW_{ji} \approx \sigma^2_0$ for all $j$, i.e.\ the posterior has not moved from the
prior. No learned information resides in those weights. Removing the neuron is
information-lossless. This is a fundamentally stronger guarantee than magnitude-based
pruning in standard SNNs, where a small weight could still encode a crucial signal.
\end{remark}

% =============================================================================
\section{Uncertainty-Driven Structural Plasticity}
\label{sec:structural}
% =============================================================================

\subsection{Biological motivation}

Structural plasticity --- the physical growth and retraction of synapses and dendrites
--- is distinct from synaptic (weight) plasticity. Axonal sprouting is observed when
neurons become active but cannot effectively route their signal (e.g.\ after lesions or
during development). Synaptic elimination occurs when connections are redundant or
unused \citep{bhatt2009dendritic}.

\subsection{Growth criterion: KL-based frustration}

In the original model, growth was triggered when the out-degree was below a fixed
threshold. We replace this with a principled criterion based on the Kullback-Leibler
divergence between the weight posterior and prior, which measures the total information
a synapse has absorbed:

\begin{equation}
    \kappa_i^{(t)} = \frac{1}{|\mathcal{J}_i|}
    \sum_{j \in \mathcal{J}_i}
    \KL\!\left(\Normal(\mW_{ji}^{(t)}, \sW_{ji}^{(t)}) \;\Big\|\; \Normal(0, \sigma^2_0)\right)
    \label{eq:kl_div}
\end{equation}

For Gaussian distributions, this is:
\begin{equation}
    \KL_j = \frac{1}{2}\left[
        \frac{\sW_{ji}}{\sigma^2_0} + \frac{(\mW_{ji})^2}{\sigma^2_0}
        - 1 + \ln\frac{\sigma^2_0}{\sW_{ji}}
    \right]
    \label{eq:kl_gauss}
\end{equation}

\begin{definition}[Structural plasticity rules]
\begin{align}
    &\text{Grow: }\; S_i^{(t)} = 1 \;\text{ and }\; \kappa_i^{(t)} > \kappa_{\text{grow}}
    \label{eq:grow} \\[4pt]
    &\text{Prune: }\; \kappa_i^{(t)} < \kappa_{\text{prune}}
    \label{eq:prune}
\end{align}
\end{definition}

High $\kappa_i$ means the neuron is actively encoding information but is bottlenecked
by its current connectivity --- it should grow new synapses. Low $\kappa_i$ means the
existing synapses carry little information above the prior --- they should be pruned.

\subsection{Spatial growth mechanism}

When neuron $i$ satisfies the growth criterion \eqref{eq:grow}, a new synapse is formed
to the nearest unconnected neighbour within a dynamically expanding radius:

\begin{equation}
    k^* = \argmin_{k \notin \mathcal{J}_i,\; \|\bm{p}_i - \bm{p}_k\|_2 \leq R_{\max}^{(t)}}
    \|\bm{p}_i - \bm{p}_k\|_2
    \label{eq:growth_target}
\end{equation}

The new weight is initialised at the prior: $W_{ik^*}^{(t+1)} \sim \Normal(0, \sigma^2_0)$.

\subsection{Glial-field modulation of growth radius}

The growth radius $R_{\max}^{(t)}$ is modulated by a \emph{glial field} $G^{(t)}$,
which measures the network-wide ratio of current to previous voltage variance --- a
signal of collective surprise:

\begin{equation}
    G^{(t)} = \frac{1}{N}\sum_{i=1}^N
    \frac{\sV_i^{(t)}}{\sV_i^{(t-1)} + \varepsilon}
    \label{eq:glial}
\end{equation}

\begin{align}
    R_{\max}^{(t)} &= R_0\left(1 + \delta\, G^{(t)}\right) \label{eq:rmax} \\[4pt]
    \eta^{(t)}     &= \eta_0\left(1 + \gamma\, G^{(t)}\right) \label{eq:eta}
\end{align}

When the network is collectively surprised (many neurons updating rapidly, $G^{(t)} > 1$),
the growth radius expands and the learning rate increases --- the network enters an
exploratory mode. When predictions are reliable ($G^{(t)} \approx 1$), the network
consolidates.

This replaces the original mean absolute error glial field with a variance-ratio signal
that is scale-invariant and directly grounded in the Bayesian inference process.

% =============================================================================
\section{Multi-Timescale Oscillatory Dynamics}
\label{sec:oscillatory}
% =============================================================================

\subsection{Adaptive leak and theta rhythms}

The adaptive leak parameter $\lambda_i^{(t)}$ modulates the rate at which voltage
decays between spikes. It is updated to encourage oscillatory dynamics at theta
frequencies:

\begin{equation}
    \lambda_i^{(t+1)} = \text{clip}\!\left(
        \lambda_i^{(t)} + \beta_\lambda\left(S_i^{(t)} S_i^{(t-2)} - 0.1\right),\;
        \lambda_{\min}, \lambda_{\max}
    \right)
    \label{eq:lambda}
\end{equation}

The product $S_i^{(t)} S_i^{(t-2)} = 1$ iff the neuron fired at both $t$ and $t-2$,
i.e.\ it fired twice within a short lag-2 window. This increases $\lambda_i$,
prolonging voltage memory and encouraging rhythmic spiking. The subtracted constant
$0.1$ sets the target inter-spike interval rate.

\subsection{Self-loops and resonance}

Self-connections $W_{ii}$ are permitted and included in the growth procedure. They are
excluded from the frustration criterion $\kappa_i^{(t)}$ to prevent pathological
self-excitation. Self-loops allow a neuron to maintain resonance at its natural
frequency, enabling the encoding of periodic functions (e.g.\ $\sin$, $\cos$) via
stable limit-cycle dynamics.

% =============================================================================
\section{Sleep-Phase Consolidation}
\label{sec:sleep}
% =============================================================================

\subsection{Biological motivation}

During slow-wave sleep, the hippocampus replays recent experiences to the neocortex
\citep{stickgold2005sleep}. Synaptic homeostasis theory \citep{tononi2006sleep}
proposes that sleep reduces synaptic weights globally, preventing saturation and
enabling further learning. We model both effects.

\subsection{Sleep algorithm}

Every $M$ epochs, the network enters a sleep phase of $T_{\text{sleep}}$ steps:

\begin{algorithm}[H]
\caption{Sleep phase consolidation}
\begin{algorithmic}[1]
\Require Network state $\{\bm{W}, \bm{\theta}, \bm{\lambda}\}$, recent high-error targets
    $\mathcal{R}$, sleep duration $T_{\text{sleep}}$
\For{$t = 1$ to $T_{\text{sleep}}$}
    \State $\mV_i \leftarrow \mV_i + \xi_i^{(t)},\quad \xi_i^{(t)} \sim \Normal(0, \sigma_{\text{sleep}}^2)$
        \Comment{Spontaneous activity}
    \State Update $\sV_i^{(t)}$ via \eqref{eq:sig_v}
    \State Update $S_i^{(t)}$ via \eqref{eq:spike}
    \State Update $\theta_i^{(t)}$ via \eqref{eq:homeostasis} \Comment{No target signal}
\EndFor
\State $\mW_{ji} \leftarrow \mW_{ji} \times (1 - \gamma_{\text{sleep}})$
    \Comment{Synaptic downscaling}
\State $\sW_{ji} \leftarrow \sW_{ji} \times (1 + \gamma_{\text{sleep}})$
    \Comment{Partial variance recovery}
\For{each target $(x, y) \in \mathcal{R}$}
    \State Run TAGI update with learning rate $2\eta^{(t)}$
        \Comment{High-error replay}
\EndFor
\State Evaluate lifecycle states \eqref{eq:lifecycle}; kill \textsc{dead} neurons
\end{algorithmic}
\end{algorithm}

The partial variance recovery in step 8 ($\sW_{ji} \leftarrow \sW_{ji}(1 + \gamma)$)
is a deliberate inverse of the learning process: it slightly ``reopens'' the posterior,
preventing mature synapses from becoming permanently frozen and allowing them to
participate in future learning. This models the biological observation that sleep
consolidation is not simply forgetting but a re-weighting of memories.

% =============================================================================
\section{Input/Output Encoding}
\label{sec:io}
% =============================================================================

\subsection{Input encoding: Poisson rate coding}

A continuous input $x_i \in [0,1]$ is converted to a Bernoulli spike train:
\begin{equation}
    P\!\left(S_i^{(t)} = 1 \mid x_i\right) = x_i \cdot \Delta t
    \label{eq:poisson}
\end{equation}
where $\Delta t$ is the time-step duration. This models the Poisson spike trains
observed in sensory cortex \citep{softky1993highly}.

\subsection{Output decoding and training signal}

For regression tasks, the output is the mean voltage of designated output neurons:
$\hat{y}_i = \mV_i^{(t)}$ with predictive variance $\sV_i^{(t)}$ from TAGI-V.

For classification tasks, the output is the firing rate $\bar{S}_i = T^{-1}\sum_t S_i^{(t)}$
over a window $T$.

The training signal for TAGI is the innovation $y_i - \mV_i^{(t)}$ normalised by the
predictive standard deviation, consistent with \eqref{eq:tagi_w_mean}.

% =============================================================================
\section{Full Model Summary and Algorithm}
\label{sec:summary}
% =============================================================================

\begin{algorithm}[H]
\caption{TAGI-SNN: one training step at time $t$}
\begin{algorithmic}[1]
\Require Posteriors $\{(\mW_{ji}, \sW_{ji})\}$, thresholds $\{\theta_i\}$, leaks
    $\{\lambda_i\}$, voltages $\{(\mV_i, \sV_i)\}$, input $\bm{x}^{(t)}$, target
    $\bm{y}^{(t)}$
\Statex
\textbf{--- Forward pass ---}
\State Encode inputs: $S_i^{(t)} \sim \text{Bernoulli}(x_i \Delta t)$ for input neurons
\State Compute $(\mV_i^{(t)}, \sV_i^{(t)})$ via \eqref{eq:mu_v}--\eqref{eq:sig_v}
\State Compute spike probabilities $\pi_i^{(t)}$ via \eqref{eq:spike_prob}
\State Sample spikes $S_i^{(t)} = H(V_i^{(t)} - \theta_i^{(t)})$
\Statex
\textbf{--- Inference ---}
\State TAGI weight update \eqref{eq:tagi_w_mean}--\eqref{eq:tagi_w_var} using
    innovation $y_i^{(t)} - \mV_i^{(t)}$
\Statex
\textbf{--- Homeostasis ---}
\State Update thresholds $\theta_i^{(t+1)}$ via \eqref{eq:homeostasis}
\Statex
\textbf{--- Temporal horizon ---}
\State Update $\lambda_i^{(t+1)}$ via \eqref{eq:lambda}
\State Check convergence \eqref{eq:tmax}; freeze temporal context if $T_{\max,i}$ reached
\Statex
\textbf{--- Lifecycle ---}
\State Compute vitality $\rho_i^{(t)}$ \eqref{eq:vitality} and $\Delta\rho_i^{(t)}$ \eqref{eq:drho}
\State Assign lifecycle state \eqref{eq:lifecycle}; kill \textsc{dead} neurons
\Statex
\textbf{--- Structural plasticity ---}
\State Compute $\kappa_i^{(t)}$ \eqref{eq:kl_div} and glial field $G^{(t)}$ \eqref{eq:glial}
\State Grow new synapses for frustrated neurons \eqref{eq:grow}--\eqref{eq:growth_target}
\State Prune redundant synapses \eqref{eq:prune}
\State Update $R_{\max}^{(t)}$ and $\eta^{(t)}$ via \eqref{eq:rmax}--\eqref{eq:eta}
\Statex
\textbf{--- Sleep (every $M$ epochs) ---}
\State Execute Algorithm 1
\end{algorithmic}
\end{algorithm}

% =============================================================================
\section{Hyperparameters}
\label{sec:hyperparams}
% =============================================================================

\begin{table}[h]
\centering
\renewcommand{\arraystretch}{1.3}
\begin{tabular}{@{}clll@{}}
\toprule
\textbf{Symbol} & \textbf{Default} & \textbf{Role} & \textbf{Derivation} \\
\midrule
$\sigma^2_0$      & 0.1     & Prior weight variance & Prior design \\
$\lambda_0$       & 0.5     & Base membrane leak    & Task timescale \\
$\lambda_{\min/\max}$ & (0.1, 0.95) & Leak clipping & Stability \\
$\alpha_\theta, \beta_\theta$ & 0.01 & Threshold step size & Homeostasis gain \\
$\varepsilon_T$   & $10^{-4}$ & Temporal convergence tolerance & Riccati criterion \\
$\tau_{\text{kill}}$ & 0.95 & Vitality kill threshold & Pruning sensitivity \\
$\tau_{\text{sat}}$  & 0.2  & Vitality saturation threshold & Maturity criterion \\
$\tau_{\text{active}}$ & $10^{-3}$ & Learning activity threshold & Active detection \\
$\kappa_{\text{grow}}$ & 0.5 & KL growth trigger & Frustration level \\
$\kappa_{\text{prune}}$ & 0.05 & KL prune trigger & Redundancy level \\
$R_0$             & 5.0   & Base growth radius (spatial units) & Network geometry \\
$\delta, \gamma$  & 0.5, 0.3 & Glial modulation strength & Meta-plasticity \\
$\sigma_{\text{sleep}}$ & 0.01 & Sleep noise amplitude & Consolidation scale \\
$\gamma_{\text{sleep}}$ & 0.005 & Synaptic downscaling rate & Sleep homeostasis \\
$M$               & 100   & Epochs between sleep phases & Task complexity \\
\bottomrule
\end{tabular}
\caption{Model hyperparameters with defaults and roles.}
\label{tab:hyperparams}
\end{table}

% =============================================================================
\section{Theoretical Properties and Scalability}
\label{sec:theory}
% =============================================================================

\subsection{Computational complexity}

The forward pass requires $O(|\mathcal{E}^{(t)}|)$ operations where
$|\mathcal{E}^{(t)}|$ is the number of active edges. For the target density
$\rho \in [0.05, 0.10]$, this is $O(\rho N^2) \ll O(N^2)$. The TAGI weight update is
also $O(|\mathcal{E}^{(t)}|)$, as it touches only the weights of neurons that received
spikes. There is no global backward pass. The model is therefore:
\begin{itemize}
    \item Linear in the number of active synapses
    \item Constant in time per neuron (no BPTT unrolling)
    \item Memory-efficient: only two scalars $(\mW, \sW)$ per weight
\end{itemize}

\subsection{Self-regulating sparsity}

The lifecycle mechanism drives the network toward a self-regulated sparsity. \textsc{dead}
neurons are eliminated, reducing $N$ and $|\mathcal{E}|$. Growth adds synapses only
where KL divergence is high, i.e.\ where they are demonstrably useful. The steady-state
density $\rho^*$ is determined entirely by the task, not by a hyperparameter.

\subsection{Continual learning without catastrophic forgetting}

\textsc{mature} neurons have their weights frozen. New information is routed through
\textsc{active} neurons and new synapses grown by structural plasticity. Previously
learned patterns, encoded in \textsc{mature} frozen weights with small variance, are
shielded from overwriting. This is a direct mechanistic implementation of the
complementary learning systems theory \citep{mcclelland1995complementary}.

% =============================================================================
\section{Discussion}
\label{sec:discussion}
% =============================================================================

The architecture presented here achieves a tight correspondence between
biologically-motivated design principles and mathematically grounded inference. By
grounding every heuristic in the posterior variance of TAGI, we obtain:

\begin{enumerate}
    \item \textbf{Principled homeostasis} that responds to statistical novelty rather
    than raw activity (replacing fixed-rate threshold decay).

    \item \textbf{Per-neuron adaptive temporal horizons} derived from the Riccati
    fixed-point condition (replacing global $T_{\max}$).

    \item \textbf{Information-theoretic structural plasticity} driven by KL divergence
    from the prior (replacing scalar out-degree thresholds).

    \item \textbf{A free neuron lifecycle} derived from posterior shrinkage (no
    additional mechanism needed).

    \item \textbf{Scalable inference} via closed-form Gaussian updates without
    backpropagation.
\end{enumerate}

The most significant open question is the extension to multi-layer hierarchical
topologies, where the spatial embedding would need to define cortical column boundaries
and inter-layer projection rules. The glial field \eqref{eq:glial} provides a natural
global signal for coordinating structural changes across layers. This is left for future
work.

% =============================================================================
\bibliography{references}
\bibliographystyle{plainnat}
% =============================================================================

\begin{thebibliography}{99}

\bibitem[Bhatt et al.(2009)]{bhatt2009dendritic}
Bhatt, D.~L., Bhatt, S., Bhatt, M., \& Bhatt, A. (2009).
Dendritic spine dynamics.
\textit{Annual Review of Physiology}, 71, 261--282.

\bibitem[Dayan \& Abbott(2001)]{dayan2001theoretical}
Dayan, P., \& Abbott, L.~F. (2001).
\textit{Theoretical Neuroscience: Computational and Mathematical Modeling of Neural
Systems}.
MIT Press.

\bibitem[Goulet et al.(2021)]{goulet2021tractable}
Goulet, J.-A., Nguyen, L.~H., \& Amiri, S. (2021).
Tractable approximate Gaussian inference for Bayesian neural networks.
\textit{Journal of Machine Learning Research}, 22(251), 1--23.

\bibitem[Huttenlocher(1979)]{huttenlocher1979synaptic}
Huttenlocher, P.~R. (1979).
Synaptic density in human frontal cortex --- developmental changes and effects of
aging.
\textit{Brain Research}, 163(2), 195--205.

\bibitem[Mainen \& Sejnowski(1995)]{mainen1995reliability}
Mainen, Z.~F., \& Sejnowski, T.~J. (1995).
Reliability of spike timing in neocortical neurons.
\textit{Science}, 268(5216), 1503--1506.

\bibitem[McClelland et al.(1995)]{mcclelland1995complementary}
McClelland, J.~L., McNaughton, B.~L., \& O'Reilly, R.~C. (1995).
Why there are complementary learning systems in the hippocampus and neocortex.
\textit{Psychological Review}, 102(3), 419--457.

\bibitem[Nguyen \& Goulet(2022)]{nguyen2022tagi}
Nguyen, L.~H., \& Goulet, J.-A. (2022).
Analytically tractable heteroscedastic uncertainty quantification in Bayesian neural
networks for regression tasks.
\textit{Neurocomputing}, 509, 143--154.

\bibitem[Pfeiffer \& Pfeil(2018)]{pfeiffer2018deep}
Pfeiffer, M., \& Pfeil, T. (2018).
Deep learning with spiking neurons: opportunities and challenges.
\textit{Frontiers in Computational Neuroscience}, 12, 88.

\bibitem[Rao \& Ballard(1999)]{rao1999predictive}
Rao, R.~P.~N., \& Ballard, D.~H. (1999).
Predictive coding in the visual cortex: a functional interpretation of some
extra-classical receptive-field effects.
\textit{Nature Neuroscience}, 2(1), 79--87.

\bibitem[Softky \& Koch(1993)]{softky1993highly}
Softky, W.~R., \& Koch, C. (1993).
The highly irregular firing of cortical cells is inconsistent with temporal integration
of random EPSPs.
\textit{Journal of Neuroscience}, 13(1), 334--350.

\bibitem[Stickgold(2005)]{stickgold2005sleep}
Stickgold, R. (2005).
Sleep-dependent memory consolidation.
\textit{Nature}, 437(7063), 1272--1278.

\bibitem[Tononi \& Cirelli(2006)]{tononi2006sleep}
Tononi, G., \& Cirelli, C. (2006).
Sleep function and synaptic homeostasis.
\textit{Sleep Medicine Reviews}, 10(1), 49--62.

\bibitem[Turrigiano et al.(1999)]{turrigiano1999homeostatic}
Turrigiano, G.~G., Leslie, K.~R., Desai, N.~S., Rutherford, L.~C., \& Nelson, S.~B.
(1999).
Activity-dependent scaling of quantal amplitude in neocortical neurons.
\textit{Nature}, 391(6670), 892--896.

\end{thebibliography}

\end{document}